\documentclass[11pt,a4paper]{article}

% ============================================================
% Paquetes
% ============================================================
\usepackage[utf8]{inputenc}
\usepackage[T1]{fontenc}
\usepackage[spanish]{babel}
\usepackage[margin=2.5cm,headheight=32pt]{geometry}
\usepackage[expansion=false]{microtype}
\usepackage{xcolor}
\usepackage{booktabs}
\usepackage{array}
\usepackage{longtable}
\usepackage{amsmath}
\usepackage{enumitem}
\usepackage{fancyhdr}
\usepackage{listings}
\usepackage[colorlinks=true,linkcolor=blue!50!black,urlcolor=blue!50!black]{hyperref}
\usepackage{titlesec}

% ============================================================
% Colores y estilos
% ============================================================
\definecolor{verde}{RGB}{46,125,50}
\definecolor{codebg}{RGB}{245,245,245}
\definecolor{codekeyword}{RGB}{0,0,150}
\definecolor{codestring}{RGB}{0,120,0}
\definecolor{codecomment}{RGB}{120,120,120}

\sloppy
\pagestyle{fancy}
\fancyhf{}
\fancyhead[R]{\footnotesize CI-0140\\Calidad de software\\I ciclo 2026}
\renewcommand{\headrulewidth}{0pt}
\fancyfoot[C]{\thepage}

\lstdefinestyle{code}{
  backgroundcolor=\color{codebg},
  basicstyle=\ttfamily\footnotesize,
  keywordstyle=\color{codekeyword}\bfseries,
  stringstyle=\color{codestring},
  commentstyle=\color{codecomment}\itshape,
  breaklines=true,
  frame=single,
  rulecolor=\color{gray!50},
  numbers=left,
  numberstyle=\tiny\color{gray},
  numbersep=6pt,
  showstringspaces=false,
  morekeywords={const,let,var,await,async,return,if,else,throw,new,export,
    default,function,class,import,from,null,true,false,static,try,catch,
    finally,extends,this},
  columns=fullflexible,
}
\lstset{style=code}

\newcommand{\campo}[2]{\par\noindent\textbf{#1:} #2\par\vspace{2pt}}
\newcommand{\campopar}[2]{\par\noindent\textbf{#1:}\par\begin{quote}\noindent #2\end{quote}}

\titleformat{\section}{\Large\bfseries}{\thesection}{1em}{}
\titleformat{\subsection}{\large\bfseries}{\thesubsection}{1em}{}
\titleformat{\subsubsection}{\normalsize\bfseries}{\thesubsubsection}{1em}{}

\title{\vspace{-2cm}Proyecto de Calidad de Software\\[0.3cm]\Huge\textbf{Fase 3}\\[0.3cm]\large Sistema POS de entregas para PYMEs -- Delivo}
\author{Aaron Santana Valdeloamr \\ Diego Cerdas Delgado}
\date{2 de julio de 2026}

\begin{document}

\maketitle
\thispagestyle{fancy}
\begin{center}
\textit{Repositorio: \texttt{tributocr-mobile} (aplicación móvil ``Delivo'')}
\end{center}
\vfill
\begin{center}
\small Universidad de Costa Rica -- Escuela de Ciencias de la Computación e Informática\\
CI-0140 Calidad de Software -- I ciclo 2026
\end{center}

\newpage
\tableofcontents
\newpage

% ============================================================
\section{Introducción}
% ============================================================

Continuando el plan de calidad definido en las Fases 1 y 2 para \textbf{Delivo}, en esta
iteración se seleccionaron dos casos de uso del \textit{product backlog} para programarlos y
verificarlos de forma completa:

\begin{enumerate}
  \item \textbf{PBI-003 -- Registro y Gestión de Órdenes de Pedido.}
  \item \textbf{PBI-004 -- Inteligencia Logística y Optimización de Rutas de Entrega.}
\end{enumerate}

Ambos ya contaban con una implementación funcional heredada del desarrollo incremental del
producto. El trabajo de esta fase se concentró en dotarlos de pruebas automatizadas reales
(no existían antes en el repositorio) y en ejecutar dos herramientas de análisis de código
sobre el proyecto completo. Las métricas de GQM que se reportan con evidencia real son las
del \textbf{Objetivo 3 (mantenibilidad)}, por ser el que corresponde directamente a los dos
casos de uso trabajados; los objetivos 1, 2 y 4 dependen de módulos que no se tocaron en
esta iteración y quedan pendientes de medir cuando se desarrollen.

% ============================================================
\section{Marco metodológico}
% ============================================================

\subsection{Características del proyecto}

\begin{longtable}{p{4.2cm} p{10.3cm}}
\toprule
\textbf{Atributo} & \textbf{Valor} \\
\midrule
\endhead
Nombre del producto & Delivo (repositorio \texttt{tributocr-mobile}) \\
Framework & Expo SDK \textasciitilde 56, React Native 0.85.3 \\
Lenguaje & TypeScript \textasciitilde 5.9.2, modo \texttt{strict} \\
Enrutamiento & \texttt{expo-router} (basado en archivos) \\
Mapas / ruteo & \texttt{@rnmapbox/maps} + \texttt{@mapbox/polyline} \\
Persistencia & Sin base de datos embebida; backend REST remoto consumido vía \texttt{axios} (el cliente móvil es un \emph{thin client}) \\
Monitoreo & \texttt{@sentry/react-native} y \texttt{expo-observe} \\
Tamaño del código & 153 archivos \texttt{.ts}/\texttt{.tsx}, 17\,712 líneas de código \\
Modelos de dominio & 9 interfaces en \texttt{types/*.d.ts}: \texttt{Order}, \texttt{Route}, \texttt{Client}, \texttt{Company}, \texttt{Product}, \texttt{PaymentMethod}, \texttt{User}, \texttt{FinancialSummary}, \texttt{Location} \\
\bottomrule
\end{longtable}

\subsection{Arquitectura}

El proyecto usa una arquitectura en capas organizada por dominio: pantallas y componentes
de UI (\texttt{app/}, \texttt{components/}), estado con Context API y \textit{hooks} de
dominio (\texttt{context/}, \texttt{hooks/}) que encapsulan \textit{fetch} + estado de
carga/error, y una capa de servicios (\texttt{services/}) donde cada dominio hereda de una
clase base común (\texttt{BaseApiService}) que envuelve los verbos HTTP. No es Redux ni
Clean Architecture formal, pero separa responsabilidades de forma clara y hace que los
servicios sean fácilmente sustituibles por \textit{mocks} en pruebas, como se aprovechó en
la Sección~\ref{sec:pruebas}.

\subsection{Proceso de desarrollo en esta iteración}

De las diez actividades de V\&V definidas en la Fase 2, las que se ejecutaron concretamente
sobre los dos casos de uso de esta fase fueron:

\begin{itemize}[leftmargin=1.3cm]
  \item \textbf{Análisis estático y \textit{code review}} (Desarrollador Principal): se
  ejecutó ESLint sobre las 153 archivos del proyecto -- resultados en la
  Sección~\ref{sec:analisis-estatico}.
  \item \textbf{Pruebas unitarias del módulo \textit{core}} (Desarrollador de Software): se
  configuró por primera vez la infraestructura de pruebas (Jest) y se escribieron 35 pruebas
  para los dos casos de uso -- resultados en la Sección~\ref{sec:pruebas}.
\end{itemize}

Al ser un equipo de dos personas, la revisión de pares se hizo directamente sobre el código
y las pruebas producidas, cada una trazable a un criterio de aceptación por su
identificador (\texttt{CP-ORD-XX} / \texttt{CP-RUT-XX}). Esta revisión tuvo un efecto
concreto: escribir las pruebas de rutas encontró un defecto real en producción, documentado
en la Sección~\ref{sec:hallazgo-bug}.

% ============================================================
\section{Calidad general del producto}
% ============================================================

\subsection{Repositorio y versionamiento}

\begin{itemize}[leftmargin=1.3cm]
  \item \textbf{Ramas:} \texttt{main}, \texttt{development}, \texttt{production}, con flujo
  \emph{feature branch} $\rightarrow$ \textit{Pull Request} $\rightarrow$ \textit{merge}.
  \item \textbf{Convención de commits:} Conventional Commits en la mayoría de los 126
  commits de \texttt{development} (72 \texttt{feat}, 6 \texttt{fix}, 3 \texttt{chore}, 1
  \texttt{style}, más los \textit{merges} de PR).
  \item \textbf{Contribuciones:} 124 commits de Aaron Santana (desarrollo incremental del
  producto en fases previas) y 2 commits de Diego Cerdas correspondientes al trabajo de esta
  iteración (infraestructura de pruebas y suite de pruebas con la corrección del defecto
  encontrado).
  \item \texttt{coverage/} se agregó a \texttt{.gitignore} para no versionar artefactos
  generados por Jest.
\end{itemize}

\subsection{Casos de uso seleccionados}

\subsubsection{PBI-003: Registro y Gestión de Órdenes de Pedido}

\campopar{Descripción}{Como Repartidor o Administrador, quiero registrar nuevas órdenes
ingresando el cliente, producto, cantidad y dirección, para mantener el control digitalizado
de los pedidos pendientes de entrega.}

\textbf{Implementación:} pantalla \texttt{app/orders/create.tsx}, servicio
\texttt{services/orders/Orders.service.ts} (\texttt{create}, \texttt{addItemToOrder},
\texttt{getOrderById}, \texttt{all}, entre otros), estado con \texttt{hooks/useOrders.ts} y
modelo \texttt{types/Order.d.ts} (\texttt{OrderStatus}: \texttt{PENDING}/\texttt{PAID}/\texttt{CANCELLED}).

\subsubsection{PBI-004: Inteligencia Logística y Optimización de Rutas de Entrega}

\campopar{Descripción}{Como Administrador, quiero crear, organizar y optimizar las rutas de
entrega basándome en la ubicación geográfica de los clientes, para minimizar la distancia
recorrida y eliminar la planificación empírica.}

\textbf{Implementación:} pantallas en \texttt{app/routes/}, servicio
\texttt{services/routes/Routes.service.ts} (\texttt{create}, \texttt{find},
\texttt{addPoint}, \texttt{startNavigation}, \texttt{completeDelivery}), estado con
\texttt{context/RouteContext.tsx} (ordena las paradas y rastrea la entrega activa) y modelo
\texttt{types/Route.d.ts}. La optimización de la secuencia de paradas y la polilínea se
calculan en el backend; el cliente móvil solo decodifica y renderiza el resultado.

\subsection{Casos de prueba}
\label{sec:casos-prueba}

Se automatizaron 22 casos de prueba entre los dos casos de uso; los más representativos:

\begin{longtable}{p{1.7cm} p{4.6cm} p{4.1cm} p{4.1cm}}
\toprule
\textbf{ID} & \textbf{Descripción} & \textbf{Entrada} & \textbf{Resultado esperado} \\
\midrule
\endhead
CP-ORD-01 & Crear una orden asociada a un cliente & \texttt{clientId: 'client-1'} & Orden creada con \texttt{status = PENDING} \\
CP-ORD-08 & Buscar órdenes por cliente o número & texto \texttt{"esquina"} / \texttt{"202"} & Listado filtrado correctamente, sin distinguir mayúsculas \\
CP-RUT-01 & Crear una ruta nueva & \texttt{\{ name: 'Ruta Centro' \}} & Ruta creada con \texttt{status = CREATED} \\
CP-RUT-09 & Ordenar las paradas al obtener la ruta & \texttt{points} con índices desordenados & Puntos ordenados ascendentemente por \texttt{index} \\
CP-RUT-11 & Completar la última parada pendiente & Ruta de una parada, entrega completada & \texttt{status} pasa a \texttt{COMPLETED} \\
\bottomrule
\end{longtable}

\subsection{Pruebas automatizadas y cobertura}
\label{sec:pruebas}

El repositorio no tenía ninguna prueba automatizada antes de esta fase. Se configuró
\texttt{jest-expo 56.0.5}, \texttt{@testing-library/react-native 13.3.3} y \texttt{jest
29.7.0}, con la cobertura acotada a los cinco módulos donde vive la lógica de negocio de
los dos casos de uso (medir el 100\,\% de la aplicación, incluyendo pantallas de módulos no
tocados en esta iteración, habría diluido la métrica sin aportar evidencia real).

\begin{longtable}{p{6.5cm} p{2.2cm} p{2.2cm} p{2.2cm}}
\toprule
\textbf{Archivo} & \textbf{\% Sentencias} & \textbf{\% Funciones} & \textbf{\% Líneas} \\
\midrule
\endhead
\texttt{Orders.service.ts} & 100\,\% & 100\,\% & 100\,\% \\
\texttt{Routes.service.ts} & 100\,\% & 100\,\% & 100\,\% \\
\texttt{useOrders.ts} & 96.29\,\% & 100\,\% & 100\,\% \\
\texttt{useRoutes.ts} & 95.00\,\% & 100\,\% & 100\,\% \\
\texttt{RouteContext.tsx} & 93.82\,\% & 100\,\% & 98.59\,\% \\
\midrule
\textbf{Total (392 LOC, 35 pruebas)} & \textbf{95.42\,\%} & \textbf{100\,\%} & \textbf{99.28\,\%} \\
\bottomrule
\end{longtable}

La cobertura por sentencias, funciones y líneas supera el 90\,\% exigido (la cobertura de
ramas quedó en 70.45\,\%, menor, porque algunas combinaciones de manejo de error se
priorizaron menos frente al tiempo disponible).

\subsubsection{Hallazgo: defecto encontrado y corregido}
\label{sec:hallazgo-bug}

Al escribir las pruebas de \texttt{RouteContext.tsx} se detectó que la ruta nunca se
marcaba como \texttt{COMPLETED} cuando la primera parada pendiente estaba en el índice
\texttt{0} (el caso más común, toda ruta recién creada). La causa: usar \texttt{||} para un
valor que puede ser legítimamente \texttt{0}, que en JavaScript es \textit{falsy}.

\begin{lstlisting}[caption={\texttt{context/RouteContext.tsx}: antes y después de la corrección}]
// Antes: "findIndex(...) || null" convierte el indice 0 en null
setCurrentPointIndex(orderedPoints.findIndex(p => p.status === 'CREATED') || null);

// Despues: comparacion explicita en vez de coercion implicita
const pendingIndex = orderedPoints.findIndex(p => p.status === 'CREATED');
setCurrentPointIndex(pendingIndex === -1 ? null : pendingIndex);
\end{lstlisting}

Sin esta corrección, un repartidor que completaba la última parada de una ruta nunca veía la
ruta transicionar a \texttt{COMPLETED}. Es la evidencia más concreta de que las pruebas
automatizadas de esta iteración tuvieron valor real más allá de una cifra de cobertura.

\subsection{Métricas de GQM ejecutadas}

Del marco GQM de la Fase 2, el Objetivo 3 (mantenibilidad) es el medible con la evidencia de
esta iteración:

\begin{longtable}{p{7.5cm} p{2.2cm} p{2.5cm}}
\toprule
\textbf{Métrica M3.2 -- Complejidad Ciclomática} & \textbf{Valor} & \textbf{Rango (Fase 2)} \\
\midrule
\endhead
Promedio (20 funciones medidas con ESLint, regla \texttt{complexity}) & 3.05 & \textcolor{verde}{Óptimo (1--10)} \\
Máxima & 5 & \textcolor{verde}{Óptimo} \\
\bottomrule
\end{longtable}

El objetivo se cumple para el código de esta iteración: ninguna función supera 5, muy por
debajo del umbral crítico de 20. La métrica M3.1 (Índice de Mantenibilidad) no pudo
calcularse: requiere el Volumen de Halstead, y las herramientas disponibles para eso
(\texttt{complexity-report}, \texttt{escomplex}) no soportan \texttt{async}/\texttt{await} ni
sintaxis moderna de TypeScript --- se documenta como limitación real en vez de reportar un
número sin respaldo. Los objetivos 1, 2 y 4 (facturación, visibilidad en exteriores, datos
sensibles) dependen de módulos fuera del alcance de esta iteración y quedan pendientes de
medir.

% ============================================================
\section{Análisis estático}
\label{sec:analisis-estatico}
% ============================================================

El análisis estático examina el código fuente sin ejecutarlo, para detectar errores,
malas prácticas o patrones riesgosos antes de llegar a producción; se complementa con el
análisis dinámico, que observa el sistema en ejecución. Para esta entrega se usaron
\textbf{ESLint} (estático) y \textbf{Sentry} (monitoreo dinámico de errores en producción),
ya integrado y usado activamente en el producto real. Se incluye pese a no ser una
herramienta estática porque permite comparar dos filosofías de calidad complementarias, como
se muestra en la Sección~\ref{sec:comparacion}.

\subsection{ESLint}

Creado por Nicholas C.\ Zakas en 2013, ESLint es el \textit{linter} estándar del ecosistema
JavaScript/TypeScript: construye un árbol de sintaxis (AST) y lo evalúa contra reglas
configurables de estilo, buenas prácticas y errores probables. El proyecto usa
\texttt{eslint 9.25.0} (\textit{flat config}) con el preset \texttt{eslint-config-expo
10.0.0}, ejecutado con \texttt{npm run lint}.

\subsection{Sentry}

Originada en 2008 como herramienta interna de Disqus y luego liberada como proyecto de
código abierto, Sentry captura excepciones y \textit{crashes} de una aplicación en
ejecución, con contexto de dispositivo/usuario y repetición visual de la sesión
(\textit{session replay}). El proyecto usa \texttt{@sentry/react-native \textasciitilde
7.11.0}, inicializado en \texttt{app/\_layout.tsx}:

\begin{lstlisting}[caption={Inicialización de Sentry (\texttt{app/\_layout.tsx})}]
Sentry.init({
  dsn: 'https://...@o4511549122084864.ingest.us.sentry.io/...',
  replaysSessionSampleRate: 0.1,
  replaysOnErrorSampleRate: 1,
  integrations: [Sentry.mobileReplayIntegration(), Sentry.feedbackIntegration()],
});
export default Sentry.wrap(ObserveRoot.wrap(RootLayout));
\end{lstlisting}

Además del \textit{wrap} global, \texttt{components/feedback/ErrorBoundary.tsx} captura
explícitamente los errores de renderizado de React vía \texttt{Sentry.captureException}.

\subsection{Resultados}

\begin{longtable}{p{6.5cm} p{3cm} p{3cm}}
\toprule
\textbf{Regla de ESLint} & \textbf{Categoría} & \textbf{Ocurrencias} \\
\midrule
\endhead
\texttt{@typescript-eslint/no-unused-vars} & Código muerto & 35 \\
\texttt{react-hooks/exhaustive-deps} & \textit{Hooks} & 17 \\
\texttt{@typescript-eslint/array-type} & Estilo & 2 \\
\texttt{import/no-named-as-default-member} & Importaciones & 2 \\
\midrule
\textbf{Total} & \textbf{0 errores, 56 advertencias} & (165 archivos, 37 con alguna advertencia) \\
\bottomrule
\end{longtable}

Cero errores en todo el proyecto es una señal fuerte de calidad; las advertencias son deuda
técnica localizable, sin bloquear el \textit{build}.

Sobre Sentry: solo dos archivos de todo el proyecto lo invocan (\texttt{app/\_layout.tsx} y
\texttt{ErrorBoundary.tsx}). Ningún \texttt{catch} de \texttt{services/} o \texttt{hooks/}
--- incluyendo los de Órdenes y Rutas --- reporta a Sentry; todos usan solo
\texttt{console.error}. En producción, un fallo de red al crear una orden no genera ningún
evento en Sentry, solo un \textit{toast} para el usuario.

\subsection{Comparación}
\label{sec:comparacion}

\begin{longtable}{p{3cm} p{5.7cm} p{5.7cm}}
\toprule
 & \textbf{ESLint} & \textbf{Sentry} \\
\midrule
\endhead
Tipo & Estático, sin ejecución & Dinámico, en producción \\
Cuándo detecta & En cada \textit{commit}/PR & Solo si el código defectuoso se ejecuta y falla \\
Qué detecta & Código muerto, malas prácticas, complejidad & Excepciones y \textit{crashes} reales, con contexto de sesión \\
Falsos negativos & No detectó el bug de la Sección~\ref{sec:hallazgo-bug} (sintaxis válida) & Ciego a errores que solo se registran con \texttt{console.error} \\
Costo & Gratuito, sin infraestructura & Cuota de eventos/\textit{replays} en el plan pagado \\
\bottomrule
\end{longtable}

Ambas herramientas son factibles en proyectos grandes: ESLint escala sin costo adicional y
es el estándar de facto del ecosistema; Sentry gana valor con el tráfico real, pero su
utilidad depende de que el equipo instrumente el código de negocio, no solo el nivel
superior de la aplicación.

% ============================================================
\section{Conclusiones}
% ============================================================

La calidad de los dos casos de uso trabajados se sustenta en evidencia verificable: pruebas
automatizadas reales que antes no existían (35 casos, cobertura de sentencias/funciones/líneas
sobre el 90\,\% exigido), análisis estático limpio (0 errores de ESLint, complejidad
ciclomática baja) y un defecto funcional real detectado y corregido gracias a las pruebas.

\textbf{Lecciones aprendidas:}

\begin{itemize}[leftmargin=1.3cm]
  \item El análisis estático y las pruebas dinámicas no son intercambiables: el bug de
  \texttt{RouteContext.tsx} era sintácticamente válido y solo una prueba de comportamiento
  lo reveló.
  \item Adoptar pruebas en un proyecto que nunca las tuvo tiene un costo de arranque
  (versión correcta de \texttt{jest-expo}, resolución de dependencias nativas), pero la capa
  de servicios y \textit{hooks}, al estar desacoplada de la UI, fue rápida de probar.
  \item Tener Sentry integrado no garantiza observabilidad si no se instrumenta en todas las
  capas; instrumentar los \texttt{catch} de la capa de servicios queda como mejora para la
  siguiente iteración.
  \item Medir solo lo que hay evidencia real para medir --- en vez de reportar los cuatro
  objetivos de GQM cuando solo dos casos de uso se trabajaron --- es en sí una práctica de
  calidad.
\end{itemize}

\end{document}
